\section{Introduction}
\label{ch:intro}

\subsection{Overview}

Every passenger car built since the mid-1990s carries a diagnostic socket
under its dashboard. Through that socket --- the OBD-II port --- the engine
control unit (ECU) will answer questions about itself once per second:
engine speed, coolant temperature, how much fuel correction it is currently
applying, how far the throttle plate is open, and dozens more. A ten-dollar
ELM327 Bluetooth adapter turns any phone into a reader for this stream.
The data is there, in every car, all the time.

What the industry does with that stream is mostly reactive. The ECU compares
individual sensors against fixed thresholds, and when one is crossed for long
enough, it stores a diagnostic trouble code (DTC) and lights the check-engine
lamp. By that point the fault is established. The slow drift that preceded it
--- a fuel trim creeping upward for weeks as an injector clogs, an idle speed
rising a little as a vacuum line cracks --- is visible in the sensor stream
long before any threshold trips, but nothing in the standard diagnostic chain
looks for it.

We built a system that looks for it. The system ingests the 1\,Hz OBD-II
stream, summarizes each 60-second window of driving into a feature vector,
and classifies the engine's state into healthy, cold-start, or one of four
developing fault classes. It runs on a laptop, takes its data from either a
recorded CSV file or a live Bluetooth adapter, explains each prediction in
terms of the sensors that drove it, and ends every session with a plain
verdict about the engine's health. Alongside the analysis pipeline we built
an Android application, OBD-LOG, that performs the vehicle-side data
collection on real cars.

One fact shapes the entire project, and we state it here rather than bury it:
there is no public dataset of real, labeled engine faults recorded through
OBD-II at 1\,Hz. Nobody has published one. Our models are therefore trained
on \emph{synthetic} faults --- physically grounded modifications injected
into 84.5 hours of real healthy driving data. Chapter~\ref{ch:results}
quantifies exactly what that means for how the results should be read.

\subsection{Problem Statement}

Given a stream of standard OBD-II sensor readings sampled at approximately
1\,Hz from an unmodified passenger car, decide --- within about a minute of
data --- whether the engine is behaving normally, warming up from a cold
start, or exhibiting the early signature of one of four common fault classes:
an intake vacuum leak, a fuel delivery restriction, a stuck coolant
temperature sensor, or a worn throttle position sensor. The decision must
come with an explanation a mechanic can check, must transfer to vehicles the
model has never seen, and must refuse to answer rather than guess when the
target vehicle does not expose the sensors a given fault requires.

Three properties of the setting make this harder than a textbook
classification exercise. The data rate is low: at 1\,Hz, combustion-level
events are invisible, and only slow aggregate signatures survive. The
labels do not exist: healthy public data is plentiful, but labeled fault
data at this rate is not, so the training signal has to be synthesized
without violating engine physics. And the deployment target varies: two cars
expose different PID sets, idle at different speeds, and trim fuel around
different baselines, so a model fitted to one car's raw values fails on the
next car for reasons that have nothing to do with faults.

\subsection{Motivation and Application}

The gap between ``data available'' and ``warning delivered'' is the
motivation. A drivetrain fault that is caught while it is still a drifting
sensor value costs a hose clamp or a sensor swap; the same fault caught at
the DTC stage may have been degrading fuel economy for months, and caught at
the breakdown stage it costs a tow. Fleet operators already pay for
telematics systems that promise early warning, but those systems are
proprietary, and the published research behind them almost never operates on
the 1\,Hz consumer-grade interface that ordinary vehicles actually expose.

Consider the concrete case of a clogging fuel injector. Early, it shows only
as a long-term fuel trim that has crept from its normal few percent to eight
or nine, correcting for a mixture that is running slightly lean. The car
drives fine; the driver notices nothing; no code is stored. Weeks later the
trim saturates at its limit, the mixture goes lean under load, and now there
is a misfire code, a rough idle, and possibly a damaged catalytic converter
downstream. The whole trajectory was visible in one PID from the first week.
A tool that flagged the drift at week one turns a converter and an injector
into just an injector, and turns a roadside stall into a scheduled
appointment. That is the value proposition in a single fault, and three of
our four fault classes tell structurally similar stories.

The intended applications are concrete. A workshop can replay a customer's
recorded drive and get a per-fault verdict with sensor-level evidence. A
student or hobbyist can point the dashboard at a live car and watch the
classifier, the fuel trims, and the SHAP explanation update together. A
fleet, with the future-work pipeline of Section~\ref{sec:future}, could
stream many vehicles into the same scoring backend. In every case the
hardware cost is a phone and an ELM327 adapter --- perhaps fifteen dollars
of equipment reading data the car already computes and publishes for free.

There is also a methodological application that outlives this particular
system. The fault injection engine of Section~\ref{sec:injector} is a
worked answer to a question many practitioners face: how do you train a
detector when the failure data does not exist? Our answer --- inject faults
into real healthy data under hard physical constraints, document every
parameter, and measure how much the scores depend on the injector's own
settings --- is reusable well beyond these four fault classes.

\subsection{Challenges}

The project's difficulties were not where we first expected them. Put
simply, the models were the easy part; the data was not.

\begin{itemize}
  \item \textbf{A corrupted-looking corpus.} Our first audit of the public
    carOBD dataset concluded that 120 of its 129 files contained physically
    impossible sensor values and had to be discarded. The real cause turned
    out to be a single-character CSV parsing defect --- a trailing comma
    that silently shifted every column one position left. Diagnosing this
    took longer than training any model in this report, and
    Section~\ref{sec:dataset} documents it in full.
  \item \textbf{Synthesizing faults without lying about physics.} A naive
    fault injector produces textbook signatures that real ECUs never emit.
    A vacuum leak on a speed-density engine, for example, does \emph{not}
    produce the large fuel-trim swing the textbooks describe, because the
    MAP sensor sees the leak directly and the ECU compensates. Getting each
    of the four fault models to match documented ECU behavior forced us to
    discard and rewrite our first implementations.
  \item \textbf{Evaluation that does not flatter itself.} Windows cut from
    the same drive overlap by up to 50 of their 60 rows. Any split that
    mixes one session's windows across train and test leaks information and
    inflates F1 by 10--20 points. We had to lock session-level splitting in
    as a tested invariant, and we had to accept that our headline numbers
    measure self-consistency with our own injector, not field detection.
  \item \textbf{Cross-vehicle reality.} The training vehicle (Toyota Etios)
    is a speed-density car with a MAP sensor; our test vehicle (Toyota
    Yaris) is MAF-based and reports no MAP at all. The first deployment on
    the healthy Yaris produced a confident, entirely wrong fault alarm, and
    fixing that failure properly reshaped the system's reporting layer.
  \item \textbf{Throughput of consumer hardware.} A real ELM327 clone polls
    7 PIDs at 1.10\,Hz but 21 PIDs at only 0.37\,Hz. Choosing which PIDs
    are worth their cost in sample rate became a design decision with
    measured trade-offs, not a footnote.
\end{itemize}

None of these five were solved by writing more machine learning code. Each
one was a data problem, a physics problem, or a hardware problem first, and
the model only got to be as good as the data underneath it allowed. That
ordering --- fix the data, then fit the model --- turned out to be the
project's actual methodology, even though it was not the one we set out
with in Week~1.

\subsection{Objectives}

\begin{enumerate}
  \item Recover and validate the full 129-session carOBD corpus, with a
    loud, tested guard against the column-shift defect that had hidden
    93\,\% of it.
  \item Build a fault injection engine for four fault classes, in ramp and
    step modes, in which every written value is clamped to engine physics
    and every injection is reproducible from stored parameters.
  \item Design a vehicle-agnostic feature representation --- 83 features
    per 60-second window, z-scored against the source vehicle's own healthy
    baseline --- that lets a model trained on one car run on another.
  \item Train and evaluate a six-class XGBoost classifier under
    session-level and family-level cross-validation, reporting synthetic
    self-consistency scores as exactly that.
  \item Complement the classifier with a per-PID forecaster and a one-class
    anomaly detector, and report honestly where they fail (long-term fuel
    trim forecasting fails).
  \item Build the OBD-LOG Android app for real-vehicle data collection over
    Bluetooth, with tiered PID selection driven by measured poll rates.
  \item Deploy the full pipeline against real recordings from an unseen
    vehicle, and design the reporting layer so that missing sensors produce
    ``Untested'' verdicts rather than fabricated scores.
  \item Deliver a live dashboard that shows the classification, its SHAP
    explanation, physics-based severity, and an end-of-read health report.
\end{enumerate}

\subsection{Ethics}

No physical fault was induced on any vehicle at any point in this project.
Every fault our models were trained on is synthetic, and every real
recording we scored came from a healthy, unmodified car. We considered
inducing a real vacuum leak under supervision and decided against it for
this phase; the documented collection protocol in the repository specifies
how to do it safely if a future team takes that step, with a mechanic
present and abort conditions defined in advance.

All driving data we collected ourselves was recorded on public roads with
the driver's knowledge and consent, at legal speeds, with the phone mounted
and the app operated by a passenger. The carOBD corpus we trained on is
public, published by its author for research use, and credited in
Chapter~\ref{ch:literature} and the references.

The system's outputs are advisory. A classifier label on a dashboard is not
a diagnosis, and the report says so on its face: every end-of-read verdict
carries its caveats in the rendered output, including the synthetic origin
of the training data and the conditions under which the verdict is
unreliable. A qualified mechanic must confirm any suspected fault before a
repair decision is made. We consider the refusal behaviors --- ``Untested''
when a required sensor is absent, ``insufficient data'' when too few windows
are evaluable --- to be ethics features as much as engineering ones, because
the alternative is a tool that manufactures confident false alarms, and we
watched an early version of our own system do exactly that.

\subsection{Report Layout}

Chapter~\ref{ch:literature} surveys prior work on OBD-II diagnostics and
machine learning for predictive maintenance, and identifies the gap this
project fills. Chapter~\ref{ch:system} describes the proposed system:
architecture, the Android data-collection app, the dataset and its recovery,
the fault injection engine, feature engineering, the models, and the
dashboard. Chapter~\ref{ch:results} presents the evaluation methodology and
results, including the real-vehicle deployment and the failures we kept.
Chapter~\ref{ch:conclusion} concludes with limitations, stated plainly, and
a prioritized future-work list. The appendices contain the key source code
listings from both repositories and the full per-class metric tables.
